\uuid{uM1Y}
\titre{Coordonnées}
\auteur{bodin ; hanani ; mzari}
\organisation{exo7}
\licence{CC BY-NC-SA 4.0}
\datecreate{2020-01-20}
\niveau{L1}
\module{Géométrie}
\chapitre{Géométrie du plan}
\sousChapitre{Distance, norme, produit scalaire}
\difficulte{3}
\theme{coordonnées}
\exo7qcmid{lille-1.geometrie-plane.022}

\contenu{
\begin{question}
\qclassification{140.01,140.02}
\qtags{motcle=coordonnées}
Dans le plan muni d'un repère orthonormé $(O,\vec{i},\vec{j})$, on considère les points $\displaystyle A(a,b)$ et $\displaystyle B(1,1)$. Comment choisir les réels $a$ et $b$ pour que le triangle de sommets $O,A,B$ soit rectangle et isocèle en $A$ ?
\begin{answers}
    \bad{$\displaystyle a=-1$ et $\displaystyle b=-1$}
    \good{$\displaystyle a=1$ et $\displaystyle b=0$}
    \good{$\displaystyle a=0$ et $\displaystyle b=1$}
    \bad{$\displaystyle a=1$ et $\displaystyle b=-1$}
\end{answers}
\begin{explanations}
On doit avoir $\displaystyle \|\overrightarrow{OA}\|=\|\overrightarrow{AB}\|$ et $\overrightarrow{OA}\cdot\overrightarrow{AB}=0$. Ceci donne ($a=1$ et $b=0$) ou ($a=0$ et $b=1$).
\end{explanations}
\end{question}
}
